\documentclass{article}
\usepackage[T1]{fontenc}
\usepackage{lmodern}
\usepackage[hidelinks]{hyperref}
\usepackage{peasy-todo}

\newcommand{\cs}[1]{\texttt{\textbackslash#1}}

\title{The \textsf{peasy-todo} package\\\large v1.0}
\author{Pietro Brach del Prever}
\date{2026/10/14}

\begin{document}
\maketitle

\begin{abstract}
  \textsf{peasy-todo} adds numbered, highlighted to-do and done notes to a document
  and collects them in a list of to-dos. It is a modified version of
  \textsf{easy-todo} by Juan Rada-Vilela.
\end{abstract}

\listoftodos

\section{Usage}

Load the package in the preamble:
\begin{verbatim}
\usepackage[<options>]{peasy-todo}
\end{verbatim}

\subsection{Commands}

\begin{description}
  \item[\cs{todo}\{\textit{note}\}] Prints a numbered red note and adds it to the
    list of to-dos. Example: \todo{Add the missing figure.}
  \item[\cs{todoi}\{\textit{note}\}] Prints only the number and adds the note to the
    list. Example: \todoi{A note that shows only in the list.}
  \item[\cs{done}\{\textit{note}\}] Prints a numbered blue note and adds it to the
    list. Example: \done{Fix the caption.}
  \item[\cs{todoii}{[\textit{color}]}\{\textit{note}\}\{\textit{list entry}\}] The
    general form. The note and the list entry can differ. The color defaults to red.
    Example: \todoii[magenta]{Short note.}{A longer entry for the list.}
  \item[\cs{listoftodos}] Prints the list of to-dos.
\end{description}

The notes use \textsf{soul} for the highlight. Thus a note cannot contain commands
that \textsf{soul} does not accept, for example \cs{gls} from \textsf{glossaries}.

\subsection{Options}

\begin{description}
  \item[\texttt{enable}] Shows the notes and the list (default).
  \item[\texttt{disable}] Hides the notes and the list.
  \item[\texttt{obeyFinal}] Hides the notes and the list when the document class
    gets the \texttt{final} option.
  \item[\texttt{section}] Gives the list a \cs{section*} heading (default).
  \item[\texttt{chapter}] Gives the list a \cs{chapter*} heading.
\end{description}

The commands \cs{todoenable}\{\texttt{true}|\texttt{false}\},
\cs{todoobeyfinal}\{\texttt{true}|\texttt{false}\}, and
\cs{todoindexas}\{\texttt{section}|\texttt{chapter}\} set the same values in the
document body.

\subsection{Customization}

Redefine \cs{todoindextitle} (default \texttt{ToDo}) to change the list heading.
Redefine \cs{todoindexpagetitle} (default \texttt{P.}) to change the page column
heading. Redefine \cs{todocolor} (default \cs{color\{red\}}) to change the list color.

\section{Dependencies}

\textsf{color}, \textsf{tocloft}, \textsf{ifthen}, \textsf{ifdraft}, and
\textsf{soul}.

\section{License}

Copyright \copyright{} 2026 Pietro Brach del Prever. Licensed under the Apache
License, Version 2.0. Based on \textsf{easy-todo} (2014/01/01) by Juan
Rada-Vilela, also under the Apache License, Version 2.0.

\end{document}
